% !TeX root = closed-systems-from-comparison-completeness.tex
% ============================================================
\chapter{Universal Classification of Representative Enrichment}
\label{ch:locus}
% ============================================================
\begin{chapterorientation}
\orientationitem{Input}{Structural closure and quotient semantics for the physical quotient \(\Phys\).}
\orientationitem{Role}{This chapter classifies the possible ways representative-level data can be added without changing the quotient-level theory.}
\orientationitem{Output}{The representative/object locus and the morphism/transport locus, with no third enrichment locus.}
\end{chapterorientation}
\section{Introduction}
Under the standing principle of closed-world admissibility
(\cref{sp:closed-world}), this chapter develops the representative-enrichment
consequences of the quotient-descend and transport-visibility clauses
(\cref{sp:quotient,sp:transport}).
With the structural closure theorem of \cref{ch:closure} fixed, it develops
the universal enrichment classification and proves that no third
enrichment locus is admissible.
\Cref{ch:bottom,ch:rotation} complete the semantic part of the closed
comparison program.  Rectangular comparison completeness yields the
canonical factorization
\[
	U\cong X_A\times X_B=:X,
\]
the intrinsic symmetry group
\[
	G:=\Aut(U,\mathcal C)
\]
acts diagonally on \(X\), and the orbit projection
\[
	\pi:X\to \Phys:=X/G
\]
is thereby fixed by the comparison datum itself.  By canonical descent,
admissible state reports are exactly the \(G\)-invariant maps
\[
	R:X\to S,
\]
equivalently the maps factoring uniquely through \(\pi\).  The quotient
\(\Phys\) therefore contains precisely the coherent state content
visible to admissible observables.
From this point onward the problem is no longer semantic.  The quotient
already records everything that survives intrinsic descent.  Any passage
back from \(\Phys\) to \(X\) must therefore introduce data absent from
quotient semantics itself.  The possible enrichment data are exhausted by
two loci that may occur together: the object layer, consisting of the choice of
representatives in the fibres of
\[
	\pi:X\to\Phys,
\]
and the morphism layer, consisting of the transport elements of \(G\)
relating those representatives along arrows.  No third locus exists.
This is not an interpretation imposed on the formalism.  It is determined
by the categorical object governing representative reconstruction.
Once \(X\), the diagonal \(G\)-action, and the quotient map \(\pi\) are
fixed, representative reconstructions are precisely functors into the
action groupoid
\[
	X\sslash G,
\]
and a functor into \(X\sslash G\) is determined exactly by its values
on objects and on morphisms.  The object and morphism loci are
therefore the two irreducible layers of data carried by the lifting
object itself.
The chapter is organized around the following classification theorem.
\begin{theorem}[Universal classification of representative enrichment]
	\label{thm:representative-enrichment-classification}
	Under the standing principle \cref{sp:closed-world}, let
	\[
		\pi:X\to \Phys:=X/G
	\]
	be the canonical orbit projection of the diagonal action of
	\[
		G=\Aut(U,\mathcal C)
	\]
	on the canonical carrier \(X=X_A\times X_B\), and let
	\[
		\pr:X\sslash G\to \Phys^{\disc}
	\]
	be the associated component functor.  Let \(\mathcal I\) be a small
	category, and let
	\[
		\overline R:\mathcal I\to \Phys^{\disc}
	\]
	be a quotient-level protocol.
	Then every representative reconstruction of \(\overline R\) is exactly
	a functor
	\[
		T:\mathcal I\to X\sslash G
		\qquad\text{with}\qquad
		\pr\circ T=\overline R.
	\]
		Whenever such a lift exists, it is exhausted by the following two
		loci of data, which may occur together and admit no third
		independent enrichment class:
	\begin{enumerate}[label=\textup{(\roman*)}]
		\item \emph{Object locus:} for each object \(i\in\Obj(\mathcal I)\), a
		      choice of representative
		      \[
			      x_i\in X
			      \qquad\text{with}\qquad
			      \pi(x_i)=\overline R(i);
		      \]
			\item \emph{Morphism locus:} for each morphism \(f:i\to j\) in
			      \(\mathcal I\), a transport element
		      \[
			      g_f\in G
			      \qquad\text{such that}\qquad
			      g_f\cdot x_i=x_j,
		      \]
		      subject to the functorial relations
		      \[
			      g_{\id_i}=e,
			      \qquad
			      g_{f_2\circ f_1}=g_{f_2}g_{f_1}.
		      \]
	\end{enumerate}
	These two loci determine the lift uniquely, and no third independent
	enrichment datum exists.
	Equivalently, the action groupoid \(X\sslash G\) is the universal and
	exhaustive recipient of all representative enrichments of quotient
	semantics.
\end{theorem}
The theorem isolates the only possible sources of non-quotient data.
Representative choice occupies the object locus.  Transport, holonomy,
the first finite triangle obstruction, and the later loop-defect globalization occupy the morphism
locus.  The remainder of the chapter makes this statement precise.
First the action groupoid is constructed and projected to the discrete
quotient \(\Phys^{\disc}\).  Quotient protocols are then formulated as
functors into \(\Phys^{\disc}\), and representative reconstruction is
expressed as a lifting problem.  Existence is treated separately by a
connected-component criterion.  The classification theorem is then
proved by direct analysis of the data carried by a functor into
\(X\sslash G\).  Only after the classification is in hand is the action
groupoid reformulated by its universal property.  Finally, the
comparison-theoretic origin of the construction is identified:
rectangular completeness is exactly the primitive that canonically
produces the carrier \(X\), the diagonal action of \(G\), and hence the
lifting groupoid itself.
This is the point at which quotient semantics passes into transport
algebra.
% ============================================================
\section{The action groupoid}
% ============================================================
\begin{definition}[Action groupoid]
	\label{def:action-gpd}
	Let a group \(G\) act on a set \(X\).  The \emph{action groupoid}
	\[
		X\sslash G
	\]
	is the category defined as follows.
	\begin{enumerate}[label=\textup{(\arabic*)}]
		\item The objects are the elements \(x\in X\).
		\item For \(x,y\in X\), a morphism \(x\to y\) is a pair
		      \[
			      (g,x)
		      \]
		      with \(g\in G\) such that \(y=g\cdot x\).
		\item If
		      \[
			      (g,x):x\to g\cdot x
			      \qquad\text{and}\qquad
			      (h,g\cdot x):g\cdot x\to hg\cdot x,
		      \]
		      then their composite is
		      \[
			      (h,g\cdot x)\circ(g,x):=(hg,x).
		      \]
		\item The identity morphism at \(x\) is
		      \[
			      \id_x:=(e,x),
		      \]
		      where \(e\in G\) is the identity element.
	\end{enumerate}
\end{definition}
\begin{lemma}
	\label{lem:action-groupoid-groupoid}
	The category \(X\sslash G\) is a groupoid.
\end{lemma}
\begin{proof}
	Associativity follows from associativity in \(G\).  For composable
	triples,
	\[
		(k,hg\cdot x)\circ\bigl((h,g\cdot x)\circ(g,x)\bigr)
		=
		(k,hg\cdot x)\circ(hg,x)
		=
		(khg,x),
	\]
	while
	\[
		\bigl((k,hg\cdot x)\circ(h,g\cdot x)\bigr)\circ(g,x)
		=
		(kh,g\cdot x)\circ(g,x)
		=
		(khg,x).
	\]
	The identity law is immediate:
	\[
		(e,g\cdot x)\circ(g,x)=(eg,x)=(g,x),
		\qquad
		(g,x)\circ(e,x)=(ge,x)=(g,x).
	\]
	Finally, every morphism \((g,x):x\to g\cdot x\) is invertible, with
	inverse
	\[
		(g^{-1},g\cdot x):g\cdot x\to x,
	\]
	since
	\[
		(g^{-1},g\cdot x)\circ(g,x)=(e,x)=\id_x
	\]
	and
	\[
		(g,x)\circ(g^{-1},g\cdot x)=(e,g\cdot x)=\id_{g\cdot x}.
	\]
	Thus every morphism is invertible, so \(X\sslash G\) is a groupoid.
\end{proof}
\begin{lemma}[Orbit characterization]
	\label{lem:gpd-orbits}
	For \(x,y\in X\), the following are equivalent.
	\begin{enumerate}[label=\textup{(\roman*)}]
		\item \(x\) and \(y\) are isomorphic in \(X\sslash G\);
		\item there exists \(g\in G\) such that
		      \[
			      y=g\cdot x;
		      \]
		\item
		      \[
			      \pi(x)=\pi(y),
		      \]
		      where \(\pi:X\to\Phys:=X/G\) is the orbit projection.
	\end{enumerate}
\end{lemma}
\begin{proof}
	A morphism \(x\to y\) in \(X\sslash G\) is by definition a pair
	\[
		(g,x):x\to g\cdot x.
	\]
	Hence such a morphism exists if and only if \(y=g\cdot x\) for some
	\(g\in G\).  Since every morphism in \(X\sslash G\) is invertible by
	\cref{lem:action-groupoid-groupoid}, this is equivalent to \(x\) and
	\(y\) being isomorphic.  Thus \textup{(i)} and \textup{(ii)} are
	equivalent.
	Condition \textup{(ii)} says exactly that \(x\) and \(y\) lie in the
	same \(G\)-orbit.  By definition of the orbit projection, this is
	equivalent to \(\pi(x)=\pi(y)\).  Thus \textup{(ii)} and \textup{(iii)}
	are equivalent.
\end{proof}
% ============================================================
\section{The orbit quotient as a discrete base}
% ============================================================
\begin{definition}[Discrete quotient category]
	\label{def:phys-disc}
	Let
	\[
		\Phys^{\disc}
	\]
	denote the discrete category whose objects are the elements of
	\(\Phys\) and whose only morphisms are identity morphisms.
\end{definition}
\begin{definition}[Component functor]
	Define
	\[
		\pr:X\sslash G\to \Phys^{\disc}
	\]
	on objects by
	\[
		\pr(x):=\pi(x),
	\]
	and on morphisms by
	\[
		\pr(g,x):=\id_{\pi(x)}.
	\]
\end{definition}
\begin{lemma}
	The assignment \(\pr\) is a functor.
\end{lemma}
\begin{proof}
	For identities,
	\[
		\pr(\id_x)=\pr(e,x)=\id_{\pi(x)}=\id_{\pr(x)}.
	\]
	For composition, let
	\[
		(g,x):x\to g\cdot x,
		\qquad
		(h,g\cdot x):g\cdot x\to hg\cdot x.
	\]
	Then
	\[
		\pr\bigl((h,g\cdot x)\circ(g,x)\bigr)
		=
		\pr(hg,x)
		=
		\id_{\pi(x)}.
	\]
	On the other hand,
	\[
		\pr(h,g\cdot x)\circ \pr(g,x)
		=
		\id_{\pi(g\cdot x)}\circ \id_{\pi(x)}.
	\]
	Since \(\pi(g\cdot x)=\pi(x)\), the right-hand side is
	\[
		\id_{\pi(x)}.
	\]
	Thus \(\pr\) preserves identities and composition.
\end{proof}
\begin{remark}
	The functor \(\pr\) forgets precisely the representative-level data that
	quotient semantics excludes: the chosen objects in \(X\) and the
	transport labels in \(G\).  What remains is the orbit class alone.
\end{remark}
% ============================================================
\section{Quotient protocols and representative lifts}
% ============================================================
\begin{definition}[Quotient-level protocol]
	Let \(\mathcal I\) be a small category.  A
	\emph{quotient-level protocol} is a functor
	\[
		\overline R:\mathcal I\to \Phys^{\disc}.
	\]
\end{definition}
Because \(\Phys^{\disc}\) is discrete, such a functor records only orbit
classes attached to the objects of \(\mathcal I\).  It carries no
representative choices and no transport information.
\begin{definition}[Representative lift]
	A \emph{representative lift} of \(\overline R\) is a functor
	\[
		T:\mathcal I\to X\sslash G
	\]
	such that
	\[
		\pr\circ T=\overline R.
	\]
\end{definition}
Thus reconstruction of representatives from quotient-level data is a
lifting problem
\[
	\begin{array}{ccc}
		           & X\sslash G                &                \\
		T \nearrow &                           & \searrow \pr   \\
		\mathcal I & \xrightarrow{\overline R} & \Phys^{\disc}.
	\end{array}
\]
\begin{remark}
	This is the groupoid-level analogue of canonical descent.  By
	\cref{thm:descent,cor:descent},
	admissible maps out of \(X\) were exactly those descending through
	\(\pi\).  Here quotient-level data are lifted back through the canonical
	groupoid lying over \(\Phys^{\disc}\).
\end{remark}
% ============================================================
\section{Lift existence}
% ============================================================
\begin{lemma}[Lift existence criterion]
	A representative lift of
	\[
		\overline R:\mathcal I\to \Phys^{\disc}
	\]
	exists if and only if \(\overline R\) is constant on each connected
	component of \(\mathcal I\).
\end{lemma}
\begin{proof}
	Assume first that a lift
	\[
		T:\mathcal I\to X\sslash G
	\]
	exists.
	Let \(i,j\in \Obj(\mathcal I)\) lie in the same connected component of
	\(\mathcal I\).  Then there is a zig-zag of morphisms joining \(i\) to
	\(j\).  Applying \(T\) yields a zig-zag of morphisms in the groupoid
	\(X\sslash G\), so all objects appearing in the image zig-zag are
	isomorphic.  By \cref{lem:gpd-orbits}, their images under \(\pi\)
	coincide.  Hence
	\[
		\overline R(i)
		=
		\pr(T(i))
		=
		\pr(T(j))
		=
		\overline R(j).
	\]
	Therefore \(\overline R\) is constant on each connected component.
	Conversely, assume that \(\overline R\) is constant on each connected
	component of \(\mathcal I\).  For each connected component \(C\), choose
	an element
	\[
		x_C\in X
	\]
	such that
	\[
		\pi(x_C)=\overline R(i)
		\qquad
		\text{for every }i\in C.
	\]
	This is possible because \(\pi\) is surjective and \(\overline R\) is
	constant on \(C\).
	Define \(T\) on objects by
	\[
		T(i):=x_C
		\qquad
		(i\in C).
	\]
	For every morphism \(f:i\to j\) inside the component \(C\), define
	\[
		T(f):=\id_{x_C}.
	\]
	Then \(T\) preserves identities and composition because identities do.
	Moreover,
	\[
		\pr(T(i))=\pi(x_C)=\overline R(i),
	\]
	so
	\[
		\pr\circ T=\overline R.
	\]
	Thus \(T\) is a representative lift.
\end{proof}
\begin{remark}
	Because the target \(\Phys^{\disc}\) is discrete, a connected quotient
	protocol can carry only one orbit per connected component.  The
	criterion states that this obvious necessary condition is also
	sufficient.
\end{remark}
% ============================================================
\section{Exhaustion of representative enrichment}
% ============================================================
\begin{proof}[Proof of \cref{thm:representative-enrichment-classification}]
	Let
	\[
		T:\mathcal I\to X\sslash G
		\qquad\text{with}\qquad
		\pr\circ T=\overline R
	\]
	be a representative reconstruction.
	A functor into \(X\sslash G\) is determined uniquely by its values on
	objects and on morphisms.
	For each object \(i\in\Obj(\mathcal I)\), the value \(T(i)\) is an
	element of \(X\).  Since
	\[
		\pr(T(i))=\overline R(i),
	\]
	the point \(T(i)\) is a representative of the orbit class
	\(\overline R(i)\).  This is exactly the object locus.
		For each morphism \(f:i\to j\), the value \(T(f)\) is a morphism in
		\(X\sslash G\) from \(T(i)\) to \(T(j)\).  By definition of the action
		groupoid, such a morphism is of the form
		\[
			T(f)=(g_f,T(i))
		\]
		for some element \(g_f\in G\) satisfying
		\[
			g_f\cdot T(i)=T(j).
		\]
		This chosen transport label is exactly the morphism locus. Without the
		later free-action hypothesis, it need not be determined by the
		endpoints alone.
	Functoriality imposes the relations
	\[
		g_{\id_i}=e
		\qquad\text{and}\qquad
		g_{f_2\circ f_1}=g_{f_2}g_{f_1}
	\]
	for composable morphisms \(f_1,f_2\).  These are compatibility
	relations on the same transport labels; they do not introduce any
	further independent datum.
	Thus every representative reconstruction is determined completely by its
	object representatives and transport labels, and there is no third
	place at which additional representative-level structure can enter.
	This is exactly the statement of
	\cref{thm:representative-enrichment-classification}.
\end{proof}
\begin{corollary}[Exact exhaustion]
	Every representative lift is exhausted by the following two loci of
	data, which may occur together:
	\begin{enumerate}[label=\textup{(\roman*)}]
		\item representative choice on objects;
		\item transport assignment on morphisms, subject to functorial
		      compatibility.
	\end{enumerate}
	No third independent invariant datum appears at the level of the
	action groupoid.
\end{corollary}
\begin{proof}
	Immediate from
	\cref{thm:representative-enrichment-classification}.
\end{proof}
\begin{theorem}[Free-action specialization: representative lifts as pointed flat \(G\)-torsor transport on the protocol]
	\label{thm:lift-flat-torsor}
	In the local free-fibre setup, assume that the orbit fibres of
	\[
		\pi:X\to\Phys
	\]
	are free \(G\)-orbits.  Let
	\[
		\overline R:\mathcal I\to \Phys^{\disc}
	\]
	be a quotient-level protocol, and let
	\[
		T:\mathcal I\to X\sslash G
	\]
	be a representative lift.
	Then \(T\) determines, and is determined by, the following data:
	\begin{enumerate}[label=\textup{(\roman*)}]
		\item for each object \(i\in\Obj(\mathcal I)\), the orbit fibre
		      \[
			      \pi^{-1}(\overline R(i)),
		      \]
		      which is a \(G\)-torsor, together with a chosen representative
		      point
		      \[
			      x_i\in \pi^{-1}(\overline R(i));
		      \]
		\item for each morphism \(f:i\to j\), the torsor isomorphism
		      \[
			      \tau_f:\pi^{-1}(\overline R(i))\to \pi^{-1}(\overline R(j))
		      \]
		      induced by the transport label of \(T(f)\), so that
		      \[
			      \tau_f(x_i)=x_j;
		      \]
		\item these isomorphisms satisfy the flatness relations
		      \[
			      \tau_{\id_i}=\id,
			      \qquad
			      \tau_{f_2\circ f_1}=\tau_{f_2}\circ\tau_{f_1}.
		      \]
	\end{enumerate}
	Equivalently, under the free-action hypothesis, a representative lift is
	exactly a pointed flat \(G\)-torsor transport system on the protocol
	category \(\mathcal I\).
\end{theorem}
\begin{proof}
	Let \(T\) be a representative lift.
	For each object \(i\in\Obj(\mathcal I)\), the protocol specifies an
	orbit
	\[
		\overline R(i)\in \Phys.
	\]
	Its fibre
	\[
		\pi^{-1}(\overline R(i))
	\]
	is a \(G\)-torsor under the given action because, by hypothesis, the
	orbit fibre is a free \(G\)-orbit.
	For each morphism \(f:i\to j\), write
	\[
		T(f)=(g_f,T(i)).
	\]
	Then left multiplication by \(g_f\) defines a map
	\[
		\tau_f:\pi^{-1}(\overline R(i))\to \pi^{-1}(\overline R(j)),
		\qquad
		x\mapsto g_f\cdot x.
	\]
	Because \(g_f\in G\) acts bijectively on \(X\), the map \(\tau_f\) is a
	torsor isomorphism.
	Functoriality of \(T\) gives
	\[
		g_{\id_i}=e
		\qquad\text{and}\qquad
		g_{f_2\circ f_1}=g_{f_2}g_{f_1},
	\]
	hence
	\[
		\tau_{\id_i}=\id,
		\qquad
		\tau_{f_2\circ f_1}=\tau_{f_2}\circ\tau_{f_1}.
	\]
	Thus the transport system is flat.
	Conversely, suppose given the orbit fibres, chosen representatives
	\(x_i\in\pi^{-1}(\overline R(i))\), and a family of torsor
	isomorphisms \(\tau_f\) satisfying the displayed flatness relations
	and the compatibility condition
	\[
		\tau_f(x_i)=x_j
	\]
	for each morphism \(f:i\to j\).
	Because the fibres are free \(G\)-torsors, each \(\tau_f\) is given by
	left multiplication by a unique element \(g_f\in G\), characterized by
	\[
		g_f\cdot x_i=x_j.
	\]
	These elements satisfy
	\[
		g_{\id_i}=e
		\qquad\text{and}\qquad
		g_{f_2\circ f_1}=g_{f_2}g_{f_1}
	\]
	because the \(\tau_f\) do.
	Hence the assignments
	\[
		T(i):=x_i,
		\qquad
		T(f):=(g_f,x_i)
	\]
	define a functor
	\[
		T:\mathcal I\to X\sslash G
	\]
	with
	\[
		\pr\circ T=\overline R.
	\]
	Thus the pointed torsor-transport data are equivalent to a
	representative lift in the free-action setting.
\end{proof}
\begin{remark}
	The classification in
	\cref{thm:representative-enrichment-classification} already yields a
	flat transport system of orbit fibres over the protocol category.
	Under the additional free-action hypothesis of
	\cref{thm:lift-flat-torsor}, those orbit fibres are \(G\)-torsors,
	and the transport system becomes the discrete antecedent of a flat
	principal \(G\)-bundle. After geometric realization of the protocol
	category, this becomes the usual notion of a flat principal
	\(G\)-bundle. In later chapters this flat transport background first
	reappears as loop defect and triangle obstruction, and only in the
	later smooth realization chain is the stabilized obstruction read as
	curvature.
\end{remark}
% ============================================================
\section{Universal property of the action groupoid}
% ============================================================
The general classification in
\cref{thm:representative-enrichment-classification} admits a universal
reformulation.
The action groupoid is the canonical groupoid over the quotient through
which every representative reconstruction factors.
\begin{theorem}[Universal property of the action groupoid]
	Let \(\mathsf C\) be a groupoid, let
	\[
		\pi_{\mathsf C}:\mathsf C\to \Phys^{\disc}
	\]
	be a functor, and let
	\[
		\iota:X\to \Obj(\mathsf C)
	\]
	be a map satisfying
	\[
		\pi_{\mathsf C}(\iota(x))=\pi(x)
		\qquad
		\text{for all }x\in X.
	\]
	Assume moreover that for every \(g\in G\) and every \(x\in X\) there
	exists a unique morphism
	\[
		\eta_{g,x}:\iota(x)\to \iota(g\cdot x)
	\]
	in \(\mathsf C\) lying over
	\[
		\id_{\pi(x)}\in \Phys^{\disc},
	\]
	and that these morphisms satisfy
	\[
		\eta_{e,x}=\id_{\iota(x)},
		\qquad
		\eta_{hg,x}=\eta_{h,g\cdot x}\circ \eta_{g,x}.
	\]
	Then there exists a unique functor
	\[
		F:X\sslash G\to \mathsf C
	\]
	such that
	\[
		F(x)=\iota(x),
		\qquad
		F(g,x)=\eta_{g,x},
		\qquad
		\pi_{\mathsf C}\circ F=\pr.
	\]
\end{theorem}
\begin{proof}
	Define \(F\) on objects by
	\[
		F(x):=\iota(x),
	\]
	and on morphisms by
	\[
		F(g,x):=\eta_{g,x}.
	\]
	This is well defined because \((g,x)\) is a morphism
	\[
		x\to g\cdot x
	\]
	in \(X\sslash G\), while \(\eta_{g,x}\) is by hypothesis a morphism
	\[
		\iota(x)\to \iota(g\cdot x)
	\]
	in \(\mathsf C\).
	For identities,
	\[
		F(\id_x)=F(e,x)=\eta_{e,x}=\id_{\iota(x)}=\id_{F(x)}.
	\]
	For composition, let
	\[
		(g,x):x\to g\cdot x,
		\qquad
		(h,g\cdot x):g\cdot x\to hg\cdot x.
	\]
	Then
	\[
		F\bigl((h,g\cdot x)\circ(g,x)\bigr)
		=
		F(hg,x)
		=
		\eta_{hg,x}.
	\]
	By the cocycle condition,
	\[
		\eta_{hg,x}
		=
		\eta_{h,g\cdot x}\circ\eta_{g,x}
		=
		F(h,g\cdot x)\circ F(g,x).
	\]
	Hence \(F\) preserves composition.
	Next, for objects,
	\[
		(\pi_{\mathsf C}\circ F)(x)=\pi_{\mathsf C}(\iota(x))=\pi(x)=\pr(x),
	\]
	and for morphisms,
	\[
		(\pi_{\mathsf C}\circ F)(g,x)=\pi_{\mathsf C}(\eta_{g,x})=\id_{\pi(x)}=\pr(g,x).
	\]
	Thus
	\[
		\pi_{\mathsf C}\circ F=\pr.
	\]
	Uniqueness is immediate, because every object of \(X\sslash G\) is an
	element \(x\in X\), and every morphism is of the form \((g,x)\).  A
	functor with the stated properties is therefore determined to take exactly
	the values defined above.
\end{proof}
\begin{remark}
	Once \(\pi:X\to\Phys\) is fixed, the action groupoid \(X\sslash G\) is
	the canonical groupoid carrying all representative reconstructions
	compatible with the \(G\)-action.  Every other such reconstruction
	factors uniquely through it.
\end{remark}
% ============================================================
\section{Sections and the object locus}
% ============================================================
\begin{proposition}[Sections and global representative choice]
	Giving a representative choice for every orbit in \(\Phys\) is
	equivalent to specifying a section
	\[
		s:\Phys\to X
	\]
	of the orbit projection \(\pi\).
\end{proposition}
\begin{proof}
	A section \(s:\Phys\to X\) satisfies
	\[
		\pi(s(p))=p
		\qquad
		\text{for every }p\in \Phys.
	\]
	Thus \(s\) assigns to each orbit \(p\) a representative
	\[
		s(p)\in \pi^{-1}(p).
	\]
	This is exactly a global representative choice.
	Conversely, suppose one has chosen, for each \(p\in \Phys\), an element
	\[
		x_p\in \pi^{-1}(p).
	\]
	Then the assignment
	\[
		s(p):=x_p
	\]
	defines a map \(s:\Phys\to X\) satisfying
	\[
		\pi(s(p))=p.
	\]
	Hence \(s\) is a section of \(\pi\).
\end{proof}
\begin{remark}
	The object locus contains neither more nor less than preferred
	representatives in orbit fibres.
\end{remark}
% ============================================================
\section{Holonomy and the morphism locus}
% ============================================================
\begin{definition}[Transport labels and holonomy]
	Let
	\[
		T:\mathcal I\to X\sslash G
	\]
	be a representative lift.  For each morphism \(f:i\to j\), write
	\[
		T(f)=(g_f,T(i)).
	\]
	If
	\[
		\gamma=f_n\circ\cdots\circ f_1
	\]
	is a composable path in \(\mathcal I\), define
	\[
		g_\gamma:=g_{f_n}\cdots g_{f_1}\in G.
	\]
	If \(\gamma\) is a loop at \(i\), the element \(g_\gamma\) is called
	the \emph{holonomy} of \(\gamma\).
\end{definition}
\begin{lemma}[Functoriality of path labels]
	For composable paths \(\gamma_1,\gamma_2\) in \(\mathcal I\), one has
	\[
		g_{\gamma_2\circ\gamma_1}=g_{\gamma_2}g_{\gamma_1}.
	\]
\end{lemma}
\begin{proof}
	Because \(T\) is a functor,
	\[
		T(\gamma_2\circ\gamma_1)=T(\gamma_2)\circ T(\gamma_1).
	\]
	Write
	\[
		T(\gamma_1)=(g_{\gamma_1},T(i)),
		\qquad
		T(\gamma_2)=(g_{\gamma_2},T(j)).
	\]
	Then
	\[
		T(\gamma_2)\circ T(\gamma_1)
		=
		(g_{\gamma_2},T(j))\circ(g_{\gamma_1},T(i))
		=
		(g_{\gamma_2}g_{\gamma_1},T(i)).
	\]
	By uniqueness of the group label in the action groupoid, this is
	exactly
	\[
		T(\gamma_2\circ\gamma_1)=(g_{\gamma_2\circ\gamma_1},T(i)).
	\]
	Hence
	\[
		g_{\gamma_2\circ\gamma_1}=g_{\gamma_2}g_{\gamma_1}.
	\]
\end{proof}
\begin{definition}[Gauge modification]
	\label{def:locus-gauge-modification}
	Let
	\[
		T:\mathcal I\to X\sslash G
	\]
	be a representative lift, and let \((h_i)_{i\in\Obj(\mathcal I)}\) be a
	family of elements of \(G\).  The \emph{gauge modification} of \(T\) by
	\((h_i)\) is the functor
	\[
		T^h:\mathcal I\to X\sslash G
	\]
	defined on objects by
	\[
		T^h(i):=h_i\cdot T(i)
	\]
	and on morphisms by
	\[
		T^h(f):=\bigl(h_j g_f h_i^{-1},\, h_i\cdot T(i)\bigr)
		\qquad
		(f:i\to j).
	\]
\end{definition}
\begin{lemma}
	The assignment \(T\mapsto T^h\) defines a representative lift of the
	same quotient protocol \(\overline R\).
\end{lemma}
\begin{proof}
	Since \(\pi(h_i\cdot T(i))=\pi(T(i))\), the object assignments remain
	over the same orbit classes.  For a morphism \(f:i\to j\),
	\[
		(h_j g_f h_i^{-1})\cdot (h_i\cdot T(i))
		=
		h_j(g_f\cdot T(i))
		=
		h_j\cdot T(j)
		=
		T^h(j),
	\]
	so the stated arrow is well defined in \(X\sslash G\).  Identities and
	composition follow from the corresponding identities in \(G\):
	\[
		h_i e h_i^{-1}=e,
		\qquad
		(h_k g_2 h_j^{-1})(h_j g_1 h_i^{-1})=h_k(g_2g_1)h_i^{-1}.
	\]
	Thus \(T^h\) is a functor.  Finally,
	\[
		\pr(T^h(i))=\pi(T(i))=\overline R(i),
	\]
	so \(\pr\circ T^h=\overline R\).
\end{proof}
\begin{theorem}[Trivial holonomy and reduction to the object locus]
	Suppose \(\mathcal I\) is a connected groupoid and let
	\[
		T:\mathcal I\to X\sslash G
	\]
	be a representative lift of
	\[
		\overline R:\mathcal I\to \Phys^{\disc}.
	\]
	Then the following are equivalent.
	\begin{enumerate}[label=\textup{(\roman*)}]
		\item Every loop in \(\mathcal I\) has trivial holonomy.
		\item There exists a gauge-equivalent representative lift
		      \[
			      T':\mathcal I\to X\sslash G
		      \]
		      with identity morphism labels on every morphism.
		\item The morphism-level enrichment of the protocol is globally
		      trivializable, so that, up to gauge, only object-level representative
		      choice remains.
	\end{enumerate}
\end{theorem}
\begin{proof}
	\emph{\textup{(i)}\(\Rightarrow\)\textup{(ii)}.}
	Fix a base object \(i_0\in\Obj(\mathcal I)\).  For each object
	\(i\in\Obj(\mathcal I)\), choose a path
	\[
		\gamma_i:i_0\to i
	\]
	in the connected groupoid \(\mathcal I\) and define
	\[
		h_i:=g_{\gamma_i}\in G.
	\]
	Because every loop has trivial holonomy, \(h_i\) is independent of the
	chosen path.  Indeed, if \(\gamma_i'\) is another path from \(i_0\) to
	\(i\), then \((\gamma_i')^{-1}\circ\gamma_i\) is a loop at \(i_0\) in
	\(\mathcal I\), so
	\[
		e
		=
		g_{(\gamma_i')^{-1}\circ\gamma_i}
		=
		g_{\gamma_i'}^{-1}g_{\gamma_i},
	\]
	hence \(g_{\gamma_i}=g_{\gamma_i'}\).
	Consider the gauge modification \(T^{h^{-1}}\) corresponding to the
	family \(h_i^{-1}\).  For a morphism \(f:i\to j\), the new label is
	\[
		h_j^{-1}g_f h_i.
	\]
	Since \(\gamma_j\) may be chosen as \(f\circ\gamma_i\), connectedness
	of the groupoid and path-independence give
	\[
		h_j=g_f h_i.
	\]
	Therefore
	\[
		h_j^{-1}g_f h_i=e.
	\]
	Thus every morphism of \(T^{h^{-1}}\) has identity label.
	\smallskip
	\noindent\emph{\textup{(ii)}\(\Rightarrow\)\textup{(iii)}.}
	If a gauge-equivalent lift has identity labels on every morphism, then
	all transport data are trivial.  Only the representatives on objects
	remain.  Thus the morphism locus reduces to the object locus.
	\smallskip
	\noindent\emph{\textup{(iii)}\(\Rightarrow\)\textup{(i)}.}
	If the morphism-level enrichment is globally trivializable, then after
	passing to a trivialized lift every morphism has label \(e\).  Hence
	every path has label \(e\), and therefore every loop has trivial
	holonomy.  Since gauge conjugation preserves triviality of loop labels,
	the original lift also has trivial holonomy on every loop.
\end{proof}
\begin{remark}
	The point is not that a literal global section \(\Phys\to X\) must
	exist on all of \(\Phys\).  The point relevant here is sharper: under
	the same pathwise hypotheses used in the theorem above, vanishing
	holonomy is exactly the condition under which the morphism locus
	disappears, up to gauge, and the lift reduces to representative
	choice.
\end{remark}
% ============================================================
\section{Rectangular completeness as the minimal primitive}
% ============================================================
\begin{theorem}[Minimal primitive for canonical lift classification]
	\label{thm:minimal-primitive-locus}
	Rectangular comparison completeness is necessary and sufficient for the
	action-groupoid lift classification developed above to be canonically
	determined from the comparison datum \((U,\mathcal C)\).
\end{theorem}
\begin{proof}
	Under the closure clause \cref{sp:closure}, assume rectangular
	completeness.  By the classification theorem at the
	comparison layer, the canonical factor map
	\[
		\Theta:U\to X_A\times X_B
	\]
	is bijective.  Hence \(U\) is canonically identified with
	\[
		X:=X_A\times X_B.
	\]
	By the diagonal action theorem, the intrinsic symmetry group
	\[
		G=\Aut(U,\mathcal C)
	\]
	acts canonically and diagonally on \(X\), and the orbit projection
	\[
		\pi:X\to \Phys:=X/G
	\]
	is therefore canonically defined.  Consequently the action groupoid
	\[
		X\sslash G,
	\]
	the component functor
	\[
		\pr:X\sslash G\to \Phys^{\disc},
	\]
	and the representative lifting problem itself are all canonically
	determined from \((U,\mathcal C)\).
	Conversely, suppose rectangular completeness fails.  Then the canonical
	factor map \(\Theta\) is not bijective.  Therefore no canonical product
	presentation
	\[
		U\cong X_A\times X_B
	\]
	is available.  Without such a canonical carrier \(X\), there is no
	canonical diagonal action of \(\Aut(U,\mathcal C)\) on a product space,
	hence no canonical orbit projection \(\pi:X\to \Phys\), and therefore
	no canonical action groupoid \(X\sslash G\).  In particular, the notions
	of representative lift, section, transport label, and holonomy are then
	not intrinsic consequences of the primitive comparison data alone.
	Thus rectangular completeness is both sufficient and necessary for the
	universal representative-lift classification to be canonically defined.
\end{proof}
\begin{corollary}[Unique minimal primitive]
	Rectangular comparison completeness is the unique minimal primitive at
	the comparison layer determining canonical diagonal redundancy and hence the
	object--transport dichotomy.
\end{corollary}
\begin{proof}
	By \cref{thm:minimal-primitive-locus}, rectangular completeness is
	necessary and sufficient for the canonical construction of the carrier
	\(X\), the diagonal action of \(G\), the quotient map \(\pi\), and the
	action groupoid \(X\sslash G\).  All subsequent lifting data depend on
	this canonical package and are therefore available exactly under
	rectangular completeness.  Hence rectangular completeness is the unique
	minimal primitive determining the two-locus classification.
\end{proof}
% ============================================================
\section{Conclusion}
\label{sec:locus-conclusion}
% ============================================================
The structural chain proved in this chapter is
\[
	\begin{aligned}
		 & (U,\mathcal C)
		\Longrightarrow
		U\cong X_A\times X_B
		\Longrightarrow
		G\curvearrowright X
		\Longrightarrow
		\pi:X\to\Phys
		\\
		 & \Longrightarrow
		\pr:X\sslash G\to \Phys^{\disc}
		\Longrightarrow
		\text{universal lift classification}.
	\end{aligned}
\]
At the representative level, every enrichment of quotient semantics is
exhausted by the following two loci, which may occur together:
\begin{enumerate}[label=\textup{(\roman*)}]
	\item representative choice on objects, equivalently the object locus;
	\item transport and holonomy on morphisms, equivalently the morphism
	      locus.
\end{enumerate}
No third locus exists.
	This chapter therefore isolates the categorical source of the two-locus
	principle that governs the remainder of the stack.  The object locus
	reappears later as representative selection and residue.  The morphism
	locus reappears as transport, route dependence, the first finite triangle
	obstruction, the later loop-defect globalization, and finally curvature after smooth realization.  This
	classification is the exact point at which quotient semantics passes
	into transport algebra.
	Accordingly, the no-third-locus statement is not heuristic but
	complete: every admissible enrichment factor enters through at least
	one of the two classified loci, and possibly both, namely
	representative selection on objects and transport data on
	morphisms.

\Cref{ch:lift} turns this classification into explicit calculus,
constructing lifts and transport cocycles in fully concrete form.

\begin{chapterclosing}
\closingitem{Established}{Representative enrichment beyond quotient semantics is exhausted by the object locus and the morphism locus; there is no third locus.}
\closingitem{Carried forward}{The two-locus theorem becomes explicit lift and transport-cocycle calculus in the next chapter.}
\end{chapterclosing}
